\subsection{What governs an agent vs. what matters to it}
\label{sec:21.2.2}
An outcome can govern what an agent does, and it can matter to the agent. That outcomes govern isn't in dispute — agentic language models pursue objectives over long horizons and hold instructions against local pressure — and it settles nothing about whether the outcomes matter to it. You also can't establish the negative.

The evidence offered — sentimentalism as an empirical psychological hypothesis \autocite{nichols2004sentimental}, compassion as a motivational state \autocite{goetz2010compassion}, guilt predicting repair where shame predicts withdrawal \autocite{tangney2007moral} — all comes from creatures where something already matters, and what it supports is an induction, which the next section states and weighs. The counterexample would be a valuational system with authority over action that doesn't depend on any state registering as mattering.\note{Frankfurt's hierarchy of desires \autocite{frankfurt1971freedom} and Watson's valuational system \autocite{watson1975free} are accounts of this shape: what has authority over the agent's action is its own evaluation, not the strength of a felt state.} Spec-based refusal, in the next section, is the attempt to build one.
